\chapter{Order statistics and the Hill estimator}
\label{chap:hill}

This chapter develops the order statistics of a finite sample and proves the weak consistency of
the Hill estimator of the extreme value index $\gamma = 1/\alpha$ of a regularly varying tail. We
follow de Haan and Ferreira, \emph{Extreme Value Theory: An Introduction}, Springer (2006), \S3.2
[dHF], and Mason, \emph{Laws of large numbers for sums of extreme values}, Ann.\ Probab.\ 10
(1982) [Mason].

\section{Order statistics}

\begin{definition}[Descending order statistics]
  \label{def:order-stat}
  \lean{ExtremeValueTheory.orderStat, ExtremeValueTheory.orderStat_of_lt,
    ExtremeValueTheory.orderStat_of_le}
  \leanok
  For a sample $x = (x_0, \ldots, x_{n-1}) \in \bR^n$, the descending order statistics
  $X_{(1)} \ge X_{(2)} \ge \cdots \ge X_{(n)}$ are the entries of $x$ sorted decreasingly. In Lean
  they are indexed from $0$: $\texttt{orderStat}\,x\,j = X_{(j+1)}$ for $j < n$ (and $0$, a junk
  value, for $j \ge n$). They are obtained by reversing the increasing sort of the tuple $x$.
\end{definition}

\begin{definition}[Exceedance count]
  \label{def:exceed-count}
  \lean{ExtremeValueTheory.exceedCount, ExtremeValueTheory.exceedCount_eq_sum}
  \leanok
  The number of exceedances of the level $a$ by the sample $x$ is
  $N_a(x) = \# \set{ i \; | \; x_i > a } = \sum_i \indof{\set{x_i > a}}$.
\end{definition}

\begin{lemma}[Basic properties of order statistics]
  \label{lem:order-stat-basic}
  \uses{def:order-stat}
  \lean{ExtremeValueTheory.orderStat_perm, ExtremeValueTheory.sum_orderStat,
    ExtremeValueTheory.orderStat_anti}
  \leanok
  The order statistics $(X_{(1)}, \ldots, X_{(n)})$ are a permutation of the sample; in particular
  $\sum_{j=1}^n f(X_{(j)}) = \sum_i f(x_i)$ for every $f$. They are decreasing:
  $X_{(j+1)} \le X_{(i+1)}$ for $i \le j < n$.
\end{lemma}
\begin{proof}
  \leanok
  The sort is a permutation, and so is its composition with the reversal $\texttt{Fin.rev}$;
  monotonicity is the monotonicity of the sorted tuple, reversed.
\end{proof}

\begin{lemma}[Duality between order statistics and exceedance counts]
  \label{lem:order-stat-duality}
  \uses{def:order-stat, def:exceed-count, lem:order-stat-basic}
  \lean{ExtremeValueTheory.lt_orderStat_iff, ExtremeValueTheory.orderStat_le_iff}
  \leanok
  For $j < n$ and $a \in \bR$: $a < X_{(j+1)}$ if and only if $N_a(x) \ge j + 1$; equivalently
  $X_{(j+1)} \le a$ if and only if $N_a(x) \le j$.
\end{lemma}
\begin{proof}
  \leanok
  If $a < X_{(j+1)}$, then $X_{(1)}, \ldots, X_{(j+1)}$ all exceed $a$, and these are $j + 1$
  distinct entries of the sample. Conversely, if $X_{(j+1)} \le a$, then only
  $X_{(1)}, \ldots, X_{(j)}$ can exceed $a$, so $N_a(x) \le j$.
\end{proof}

\begin{lemma}[Measurability]
  \label{lem:order-stat-measurable}
  \uses{def:order-stat, def:exceed-count}
  \lean{ExtremeValueTheory.measurable_exceedCount, ExtremeValueTheory.measurable_orderStat}
  \leanok
  If $X_0, \ldots, X_{n-1}$ are measurable random variables, then the exceedance counts and the
  order statistics of $(X_0, \ldots, X_{n-1})$ are measurable.
\end{lemma}
\begin{proof}
  \leanok
  $N_a$ is a finite sum of indicators of measurable sets. For the order statistics, by
  Lemma~\ref{lem:order-stat-duality} the sets $\set{ a < X_{(j+1)} }$ are measurable for all $a$.
\end{proof}

\begin{theorem}[Distribution of exceedance counts and order statistics]
  \label{thm:exceed-count-binomial}
  \uses{lem:order-stat-duality, lem:order-stat-measurable}
  \lean{ExtremeValueTheory.measure_lt_orderStat_eq, ExtremeValueTheory.measure_exceedCount_eq,
    ExtremeValueTheory.measure_lt_orderStat_eq_sum}
  \leanok
  For random variables $X_0, \ldots, X_{n-1}$ and $j < n$,
  $\PR[X_{(j+1)} > a] = \PR[N_a \ge j + 1]$. If moreover the $X_i$ are independent with common
  exceedance probability $p = \PR[X_i > a]$, then $N_a$ is binomial,
  \begin{align*}
    \PR[N_a = m] \; = \; \binom{n}{m} p^m (1 - p)^{n - m} ,
    \qquad
    \PR[X_{(j+1)} > a] \; = \; \sum_{m = j+1}^{n} \binom{n}{m} p^m (1 - p)^{n - m} .
  \end{align*}
\end{theorem}
\begin{proof}
  \leanok
  The first identity is Lemma~\ref{lem:order-stat-duality} pointwise. The event $\set{N_a = m}$ is
  the disjoint union over $m$-element subsets $S$ of indices of the events
  $\bigcap_{i \in S} \set{X_i > a} \cap \bigcap_{i \notin S} \set{X_i \le a}$, each of probability
  $p^m (1 - p)^{n-m}$ by independence. Summing over $m \ge j + 1$ gives the last formula.
\end{proof}

\section{The Hill estimator}

\begin{definition}[Hill estimator]
  \label{def:hill}
  \uses{def:order-stat}
  \lean{ExtremeValueTheory.hill}
  \leanok
  The Hill estimator based on the top $k + 1$ order statistics of a sample $x \in \bR^n$ is
  \begin{align*}
    H_{k,n} \; = \; \frac{1}{k} \sum_{i=1}^{k} \big( \log X_{(i)} - \log X_{(k+1)} \big) .
  \end{align*}
\end{definition}

\begin{definition}[Log-excess]
  \label{def:log-excess}
  \lean{ExtremeValueTheory.logExcess, ExtremeValueTheory.logExcess_nonneg,
    ExtremeValueTheory.logExcess_of_le, ExtremeValueTheory.logExcess_of_ge,
    ExtremeValueTheory.logExcess_anti, ExtremeValueTheory.measurable_logExcess}
  \leanok
  The log-excess of $y$ over the threshold $t > 0$ is
  $\log^+(y/t) = \log\big( \max(y, t)/t \big)$. It is nonnegative and measurable in $y$, equals $0$
  for $y \le t$ and $\log y - \log t$ for $y \ge t$, and is antitone in $t$.
\end{definition}

\begin{lemma}[Log-excess representation of the Hill estimator]
  \label{lem:hill-log-excess}
  \uses{def:hill, def:log-excess, lem:order-stat-basic}
  \lean{ExtremeValueTheory.hill_eq_sum_logExcess}
  \leanok
  If $k < n$ and $X_{(k+1)} > 0$, then
  $H_{k,n} = \frac{1}{k} \sum_{i=0}^{n-1} \log^+ \big( x_i / X_{(k+1)} \big)$.
\end{lemma}
\begin{proof}
  \leanok
  By Lemma~\ref{lem:order-stat-basic}, the sum over the sample equals the sum over the order
  statistics. The terms with $j \ge k + 1$ vanish since $X_{(j)} \le X_{(k+1)}$, the term $j = k + 1$
  vanishes, and for $j \le k$ the term is $\log X_{(j)} - \log X_{(k+1)}$.
\end{proof}

\begin{lemma}[A weak law of large numbers for triangular arrays]
  \label{lem:chebyshev-wlln}
  \lean{ExtremeValueTheory.tendsto_measure_abs_sum_div_sub}
  \leanok
  Let $X_0, X_1, \ldots$ be i.i.d., $g_n$ measurable with $g_n(X_0) \in L^2$, and $b_n > 0$ with
  $n \, \EX[g_n(X_0)]/b_n \to m$ and $n \, \EX[g_n(X_0)^2]/b_n^2 \to 0$. Then
  \begin{align*}
    \frac{1}{b_n} \sum_{i<n} g_n(X_i) \; \to \; m \qquad \text{in probability.}
  \end{align*}
\end{lemma}
\begin{proof}
  \leanok
  By pairwise independence, the variance of the normalized sum is at most
  $n \, \EX[g_n(X_0)^2]/b_n^2 \to 0$, while its mean tends to $m$. Chebyshev's inequality
  concludes.
\end{proof}

\begin{definition}[Survival function of a random variable]
  \label{def:survival}
  \lean{ExtremeValueTheory.survival, ExtremeValueTheory.survival_anti,
    ExtremeValueTheory.survival_nonneg, ExtremeValueTheory.survival_pos,
    ExtremeValueTheory.survival_eq_one_sub_cdf, ExtremeValueTheory.continuousWithinAt_survival,
    ExtremeValueTheory.tendsto_survival_atTop}
  \leanok
  The survival function of a real random variable $Y$ is $\overline{F}(x) = \PR[Y > x]$. It is
  nonnegative, antitone and right-continuous, equals $1 - F_Y(x)$ for the c.d.f.\ $F_Y$ of $Y$,
  and tends to $0$ at $+\infty$. If it is regularly varying, it is everywhere positive.
\end{definition}

\begin{lemma}[Uniform exponential decay along geometric levels]
  \label{lem:survival-exp-le}
  \uses{def:survival, thm:rv-potter-monotone}
  \lean{ExtremeValueTheory.exists_survival_exp_le}
  \leanok
  If $\overline{F}$ is regularly varying with index $-\alpha$, $\alpha > 0$, there is $x_0 > 0$
  such that $\overline{F}(s e^u) \le 2 \, \overline{F}(s) \, e^{-(\alpha/2) u}$ for all $s \ge x_0$
  and $u \ge 0$.
\end{lemma}
\begin{proof}
  \leanok
  This is the Potter bound (Theorem~\ref{thm:rv-potter-monotone}) with $\delta = \alpha/2$ and
  constant $2$, applied with $c = e^u \ge 1$.
\end{proof}

\begin{lemma}[Moments of the log-excess]
  \label{lem:log-excess-moments}
  \uses{def:log-excess, def:survival, lem:survival-exp-le}
  \lean{ExtremeValueTheory.setOf_lt_logExcess, ExtremeValueTheory.integral_logExcess,
    ExtremeValueTheory.integrableOn_mul_exp_neg, ExtremeValueTheory.lintegral_logExcess_sq_le,
    ExtremeValueTheory.exists_memLp_logExcess}
  \leanok
  For $s, u > 0$, $\set{ \log^+(Y/s) > u } = \set{ Y > s e^u }$, and hence (layer cake formula)
  \begin{align*}
    \EX \big[ \log^+(Y/s) \big] \; = \; \int_0^\infty \overline{F}(s e^u) \, \ud u .
  \end{align*}
  If $\overline{F}$ is regularly varying with index $-\alpha$, $\alpha > 0$, there are $x_0, C$ such
  that for $s \ge x_0$ the log-excess $\log^+(Y/s)$ is square integrable with
  $\EX \big[ \log^+(Y/s)^2 \big] \le C \, \overline{F}(s)$.
\end{lemma}
\begin{proof}
  \leanok
  The layer cake formula gives $\EX[\log^+(Y/s)^2] = \int_0^\infty 2u \, \overline{F}(s e^u) \, \ud u$,
  and Lemma~\ref{lem:survival-exp-le} bounds it by
  $4 \overline{F}(s) \int_0^\infty u \, e^{-(\alpha/2)u} \, \ud u < \infty$.
\end{proof}

\begin{lemma}[Integrated regular variation]
  \label{lem:tail-integral-log-excess}
  \uses{lem:log-excess-moments, lem:survival-exp-le, def:regularly-varying}
  \lean{ExtremeValueTheory.integral_exp_neg_mul_Ioi, ExtremeValueTheory.tendsto_integral_logExcess_div}
  \leanok
  If $\overline{F}$ is regularly varying with index $-\alpha$, $\alpha > 0$, then
  \begin{align*}
    \lim_{s \to \infty} \frac{\EX \big[ \log^+(Y/s) \big]}{\overline{F}(s)} \; = \;
    \lim_{s \to \infty} \int_0^\infty \frac{\overline{F}(s e^u)}{\overline{F}(s)} \, \ud u
    \; = \; \int_0^\infty e^{-\alpha u} \, \ud u \; = \; \frac{1}{\alpha} .
  \end{align*}
\end{lemma}
\begin{proof}
  \leanok
  Pointwise $\overline{F}(s e^u)/\overline{F}(s) \to e^{-\alpha u}$ by regular variation, and the
  integrands are dominated by $2 e^{-(\alpha/2) u}$ (Lemma~\ref{lem:survival-exp-le}), so dominated
  convergence applies.
\end{proof}

\begin{lemma}[Intermediate thresholds]
  \label{lem:hill-threshold}
  \lean{ExtremeValueTheory.upperQuantile_spec_of_tendsto_atTop}
  \leanok
  Let $\overline{F}$ be antitone and right-continuous with $\overline{F}(x) \to 0$ as
  $x \to \infty$, and let $0 < p < \overline{F}(M)$. Then $u = \inf \set{ x \; | \; \overline{F}(x) \le p }$
  satisfies $\overline{F}(u) \le p$, $p < \overline{F}(x)$ for $x < u$, and $M \le u$.
\end{lemma}
\begin{proof}
  \leanok
  As in Lemma~\ref{lem:gumbel-quantile}.
\end{proof}

\begin{theorem}[Weak consistency of the Hill estimator]
  \label{thm:hill-consistency}
  \uses{lem:hill-log-excess, lem:chebyshev-wlln, lem:log-excess-moments,
    lem:tail-integral-log-excess, lem:hill-threshold, thm:exceed-count-binomial,
    lem:order-stat-duality, def:regularly-varying}
  \lean{ExtremeValueTheory.tendstoInMeasure_hill}
  \leanok
  Let $X_0, X_1, \ldots$ be i.i.d.\ real random variables whose survival function
  $\overline{F}(x) = \PR[X_0 > x]$ is regularly varying with index $-\alpha$, $\alpha > 0$, and let
  $k = k(n)$ be an intermediate sequence: $k(n) \to \infty$ and $k(n)/n \to 0$. Then the Hill
  estimator $H_{k(n),n}$ of $X_0, \ldots, X_{n-1}$ converges in probability to $1/\alpha$
  [dHF, Theorem~3.2.2], [Mason].
\end{theorem}
\begin{proof}
  \leanok
  Instead of the Rényi representation of exponential spacings used in [dHF], we use the tail
  empirical process argument in the spirit of [Mason], which only needs Chebyshev's inequality.

  \emph{Step 1 (deterministic).} By Lemma~\ref{lem:hill-log-excess},
  $H_{k,n} = \frac{1}{k} \sum_i \log^+(X_i / X_{(k+1)})$, and $s \mapsto \sum_i \log^+(X_i/s)$ is
  antitone (Definition~\ref{def:log-excess}).

  \emph{Step 2 (thresholds).} Choose deterministic thresholds
  $u_n = \inf \set{ x \; | \; \overline{F}(x) \le k/n }$. By Lemma~\ref{lem:hill-threshold},
  right-continuity and regular variation, $u_n \to \infty$, $n \overline{F}(u_n)/k \to 1$, and
  hence $n \overline{F}(c u_n)/k \to c^{-\alpha}$ for every $c > 0$.

  \emph{Step 3 (laws of large numbers).} Fix $c > 0$. By Lemma~\ref{lem:chebyshev-wlln} with
  $b_n = k(n)$:
  the exceedance count satisfies $N_{c u_n}/k \to c^{-\alpha}$ in probability (its summands are
  indicators with mean $\overline{F}(c u_n)$, and the second moment condition holds since
  $k \to \infty$); and $\frac{1}{k} \sum_i \log^+(X_i/(c u_n)) \to c^{-\alpha}/\alpha$ in
  probability, the mean being computed by Lemma~\ref{lem:tail-integral-log-excess} and the second
  moment being $O(\overline{F}(c u_n)) = O(k/n)$ by Lemma~\ref{lem:log-excess-moments}.

  \emph{Step 4 (sandwich).} By the duality
  $X_{(k+1)} > a \iff N_a \ge k + 1$ (Lemma~\ref{lem:order-stat-duality}), Step 3 gives
  $c_1 u_n < X_{(k+1)} \le c_2 u_n$ with probability tending to $1$ for any $c_1 < 1 < c_2$. On
  this event, monotonicity (Step 1) sandwiches $H_{k,n}$ between the averages of log-excesses
  over $c_2 u_n$ and over $c_1 u_n$, which are close to $c_2^{-\alpha}/\alpha$ and
  $c_1^{-\alpha}/\alpha$. Letting $c_1, c_2 \to 1$ concludes.
\end{proof}

\begin{corollary}[Hill estimator for Pareto samples]
  \label{cor:hill-pareto}
  \uses{thm:hill-consistency, lem:rv-examples}
  \lean{ExtremeValueTheory.tendstoInMeasure_hill_pareto}
  \leanok
  If $X_0, X_1, \ldots$ are i.i.d.\ Pareto$(\alpha)$, i.e.\ $\PR[X_0 > x] = x^{-\alpha}$ for
  $x \ge 1$, and $k(n) \to \infty$, $k(n)/n \to 0$, then $H_{k(n),n} \to 1/\alpha$ in probability.
\end{corollary}
\begin{proof}
  \leanok
  The survival function agrees with $x^{-\alpha}$ for $x \ge 1$, hence is regularly varying with
  index $-\alpha$; apply Theorem~\ref{thm:hill-consistency}.
\end{proof}

\section{The Pickands and moment estimators}

\begin{definition}[Pickands estimator]
  \label{def:pickands}
  \uses{def:order-stat}
  \lean{ExtremeValueTheory.pickands}
  \leanok
  $\hat\xi_{k,n} = (\log 2)^{-1} \log \dfrac{X_{(k+1)} - X_{(2k+1)}}{X_{(2k+1)} - X_{(4k+1)}}$
  [dHF, \S3.3], [EKM, \S6.4.2].
\end{definition}

\begin{theorem}[Consistency of the Pickands estimator]
  \label{thm:pickands-consistency}
  \uses{def:pickands, thm:pbdh-general, thm:pbdh-iff}
  \lean{ExtremeValueTheory.tendstoInMeasure_pickands, ExtremeValueTheory.tendstoInMeasure_pickands_exponential,
    ExtremeValueTheory.gpdSurvivalInv_pickands_ratio, ExtremeValueTheory.tendsto_measure_abs_exceedCount_div_sub,
    ExtremeValueTheory.tendsto_measure_not_mem_orderStat}
  \leanok
  For iid samples with c.d.f.\ $F$ in the domain of attraction of $H_\xi$, $\xi \in \bR$
  arbitrary, and $k \to \infty$, $k/n \to 0$, the Pickands estimator converges to $\xi$ in
  probability [EKM, Theorem~6.4.1], [dHF, \S3.3].
\end{theorem}
\begin{proof}
  \leanok
  Let $u_n$ be the $(1 - 8k/n)$-quantile and $\beta$ the scale from
  Pickands--Balkema--de Haan. The number of exceedances of $u_n + \beta(u_n) x$, divided by $k$,
  converges in probability to $8 \bar G_{\xi}(x)$, which localizes
  $(X_{(jk+1)} - u_n)/\beta(u_n)$ near $\bar G_\xi^{-1}(j/8)$ for $j = 1, 2, 4$. The ratio of
  differences of these three points is exactly $2^\xi$.
\end{proof}

\begin{definition}[Log-moments and the moment estimator]
  \label{def:moment-estimator}
  \uses{def:hill}
  \lean{ExtremeValueTheory.logMoment, ExtremeValueTheory.logMoment_one, ExtremeValueTheory.momentEstimator}
  \leanok
  $M^{(j)}_{k,n} = k^{-1} \sum_{i<k} (\log X_{(i+1)} - \log X_{(k+1)})^j$, so $M^{(1)}$ is the
  Hill estimator, and the moment estimator is
  $\hat\gamma = M^{(1)} + 1 - \tfrac12 \big(1 - (M^{(1)})^2/M^{(2)}\big)^{-1}$ [dHF, \S3.5].
\end{definition}

\begin{theorem}[Consistency of the log-moments and the moment estimator, $\xi > 0$]
  \label{thm:moment-consistency}
  \uses{def:moment-estimator, thm:hill-consistency}
  \lean{ExtremeValueTheory.tendstoInMeasure_logMoment, ExtremeValueTheory.tendstoInMeasure_momentEstimator,
    ExtremeValueTheory.tendstoInMeasure_momentEstimator_of_isInMDAGEV, ExtremeValueTheory.tendstoInMeasure_hill_of_isInMDAGEV,
    ExtremeValueTheory.tendstoInMeasure_momentEstimator_pareto, ExtremeValueTheory.tendsto_integral_logExcess_rpow_div}
  \leanok
  Under the hypotheses of Theorem~\ref{thm:hill-consistency},
  $M^{(j)}_{k,n} \to j!/\alpha^j$ in probability for $j \ge 1$, and the moment estimator converges
  to $1/\alpha$. Equivalently, for $F$ in the domain of attraction of $H_\xi$ with $\xi > 0$, both
  the Hill and moment estimators are consistent for $\xi$. The case $\xi \le 0$ is not
  formalized.
\end{theorem}
\begin{proof}
  \leanok
  As for the Hill estimator, with $\E[\log^+(Y/s)^p] / \bar F(s) \to \Gamma(p+1)/\alpha^p$ in place
  of the first log-moment.
\end{proof}

\section{Extreme quantiles and tail probabilities}

\begin{definition}[Weissman and tail-probability estimators]
  \label{def:weissman}
  \uses{def:hill, def:order-stat}
  \lean{ExtremeValueTheory.upperQuantile, ExtremeValueTheory.weissman, ExtremeValueTheory.tailProbEstimator, ExtremeValueTheory.IsBoundedInProbability}
  \leanok
  For $p_n \to 0$, the quantile $x_{p} = \inf\{x : \bar F(x) \le p\}$ is estimated by
  $\hat x_{p_n} = X_{(k+1)} \, (k/(n p_n))^{\hat\gamma}$, and $\bar F(y_n)$ by
  $\hat p(y_n) = (k/n)(y_n / X_{(k+1)})^{-1/\hat\gamma}$, with $\hat\gamma$ the Hill estimator
  [Weissman 1978], [dHF, \S4.3], [EKM, \S6.4.1].
\end{definition}

\begin{theorem}[Ratio consistency of the Weissman estimator]
  \label{thm:weissman}
  \uses{def:weissman, thm:hill-consistency}
  \lean{ExtremeValueTheory.tendstoInMeasure_weissman_div_of_bias, ExtremeValueTheory.tendstoInMeasure_weissman_div,
    ExtremeValueTheory.tendstoInMeasure_weissman_div_of_isBoundedInProbability,
    ExtremeValueTheory.tendstoInMeasure_weissman_div_pareto, ExtremeValueTheory.tendstoInMeasure_orderStat_div_upperQuantile,
    ExtremeValueTheory.tendsto_upperQuantile}
  \leanok
  Let $X_i$ be iid with $\bar F \in RV_{-\alpha}$, $k \to \infty$, $k/n \to 0$ and
  $d_n = k/(n p_n)$. Then $\hat x_{p_n}/x_{p_n} \to 1$ in probability in each of these cases:
  \begin{itemize}
    \item $d_n$ stays in a compact subset of $(0, \infty)$ (intermediate quantiles);
    \item $\sqrt{k}(\hat\gamma - 1/\alpha)$ is bounded in probability, $\log d_n / \sqrt k \to 0$, and the
      bias condition $x_{p_n} / (x_{k/n} d_n^{1/\alpha}) \to 1$ holds, which is automatic for exact
      Pareto tails. This covers $n p_n \to c \in [0, \infty)$, beyond the sample range.
  \end{itemize}
  [dHF, Theorem~4.3.1] (consistency part).
\end{theorem}
\begin{proof}
  \leanok
  Write $\hat x_{p_n}/x_{p_n} = (X_{(k+1)}/x_{k/n}) \cdot d_n^{\hat\gamma - 1/\alpha} \cdot
  x_{k/n} d_n^{1/\alpha} / x_{p_n}$. The first factor tends to $1$ by localization of the order
  statistic, the second because $(\hat\gamma - 1/\alpha) \log d_n \to 0$, and the third is the bias.
\end{proof}

\begin{theorem}[Ratio consistency of the tail-probability estimator]
  \label{thm:tail-prob-estimator}
  \uses{thm:weissman}
  \lean{ExtremeValueTheory.tendstoInMeasure_tailProbEstimator_div_of_bias, ExtremeValueTheory.tendstoInMeasure_tailProbEstimator_div,
    ExtremeValueTheory.tendstoInMeasure_tailProbEstimator_div_of_isBoundedInProbability}
  \leanok
  Under the same three sets of hypotheses (with $e_n = y_n/x_{k/n}$ in place of $d_n$),
  $\hat p(y_n)/\bar F(y_n) \to 1$ in probability [dHF, \S4.4].
\end{theorem}
\begin{proof}
  \leanok
  As for Theorem~\ref{thm:weissman}.
\end{proof}

\begin{definition}[Return levels]
  \label{def:return-level}
  \lean{ExtremeValueTheory.gevFun, ExtremeValueTheory.returnLevel, ExtremeValueTheory.gevFun_returnLevel}
  \leanok
  The $m$-block return level of $H_\xi$ is $z_m = ((-\log(1 - 1/m))^{-\xi} - 1)/\xi$
  ($-\log(-\log(1 - 1/m))$ for $\xi = 0$), the level with $H_\xi(z_m) = 1 - 1/m$ [Coles, \S3.1.3].
\end{definition}

\section{Upper order statistics and records}

\begin{theorem}[Limit law of the $k$-th largest]
  \label{thm:kth-largest}
  \uses{def:order-stat}
  \lean{ExtremeValueTheory.measureReal_orderStat_le_eq_sum, ExtremeValueTheory.tendsto_measureReal_orderStat_le,
    ExtremeValueTheory.tendsto_measureReal_affine_orderStat_le,
    ExtremeValueTheory.exists_tendsto_measureReal_affine_orderStat_le_of_isInMDAGumbel}
  \leanok
  For iid samples with $n(1 - F(u_n)) \to \tau \in [0, \infty)$,
  $\PR[X_{(k)} \le u_n] \to e^{-\tau} \sum_{j < k} \tau^j/j!$. Consequently, if
  $F^n(a_n x + b_n) \to H(x) \in (0, 1)$, then
  $\PR[(X_{(k)} - b_n)/a_n \le x] \to H(x) \sum_{j<k} (-\log H(x))^j/j!$
  [EKM, Theorem~4.2.3], [LLR, Theorem~2.2.2].
\end{theorem}
\begin{proof}
  \leanok
  $\{X_{(k)} \le u_n\}$ is the event that fewer than $k$ observations exceed $u_n$; the
  exceedance count is binomial with mean $\to \tau$, hence asymptotically Poisson.
\end{proof}

\begin{theorem}[Joint law of the two largest]
  \label{thm:two-largest}
  \uses{thm:kth-largest}
  \lean{ExtremeValueTheory.measureReal_orderStat_zero_le_and_orderStat_one_le,
    ExtremeValueTheory.tendsto_measureReal_orderStat_zero_le_and_orderStat_one_le,
    ExtremeValueTheory.tendsto_measureReal_affine_orderStat_zero_one}
  \leanok
  If $v_n \le u_n$, $n(1 - F(u_n)) \to \tau_1$ and $n(1 - F(v_n)) \to \tau_2$, then
  $\PR[X_{(1)} \le u_n, X_{(2)} \le v_n] \to e^{-\tau_2}(1 + \tau_2 - \tau_1)$ [EKM, \S4.2].
\end{theorem}
\begin{proof}
  \leanok
  Exactly: $\PR[X_{(1)} \le u, X_{(2)} \le v] = F(v)^n + n (F(u) - F(v)) F(v)^{n-1}$.
\end{proof}

\begin{theorem}[Records]
  \label{thm:records}
  \lean{ExtremeValueTheory.IsRecord, ExtremeValueTheory.recordCount, ExtremeValueTheory.measureReal_isRecord, ExtremeValueTheory.measureReal_isRecord_and_isRecord,
    ExtremeValueTheory.integral_recordCount, ExtremeValueTheory.integral_sq_recordCount_sub_le,
    ExtremeValueTheory.tendstoInMeasure_recordCount_div_log}
  \leanok
  For an iid sequence with atomless marginal, the $k$-th observation (counting from $1$) is a record
  with probability $1/k$, record indicators are pairwise independent, the number of records
  $N_n$ among the first $n$ observations satisfies $\E N_n = H_n = \sum_{k \le n} 1/k$ and
  $\Var N_n \le H_n$, and $N_n / \log n \to 1$ in probability [EKM, \S5.4], [Resnick, \S4.1].
  Full mutual independence of the indicators (R\'enyi) is not formalized.
\end{theorem}
\begin{proof}
  \leanok
  By exchangeability, the maximum of the first $k$ observations is equally likely to be any of
  them, and there are no ties. A second exchangeability argument gives pairwise independence;
  Chebyshev's inequality concludes.
\end{proof}

\section{Asymptotic normality of the Hill estimator}

\begin{lemma}[R\'enyi representation]
  \label{lem:renyi}
  \uses{def:order-stat}
  \lean{ExtremeValueTheory.lintegral_orderStat_excess, ExtremeValueTheory.map_orderStat_excess}
  \leanok
  For an iid standard exponential sample of size $n$ and $k < n$, the excesses
  $E_{(i)} - E_{(k+1)}$, $i \le k$, are (conditionally on $E_{(k+1)}$) iid standard exponential
  [dHF, Lemma~3.2.3].
\end{lemma}
\begin{proof}
  \leanok
  Count exceedances and split off coordinates, using the memoryless property.
\end{proof}

\begin{theorem}[Asymptotic normality of the Hill estimator]
  \label{thm:hill-normality}
  \uses{lem:renyi, thm:second-order-bound, thm:hill-consistency}
  \lean{ExtremeValueTheory.survivalQuantile, ExtremeValueTheory.tendstoInDistribution_hill_pareto,
    ExtremeValueTheory.tendstoInDistribution_hill_of_tendsto_zero, ExtremeValueTheory.tendstoInDistribution_hill,
    ExtremeValueTheory.tendstoInDistribution_meanExcess, ExtremeValueTheory.map_eq_expMeasure_map_survivalQuantile}
  \leanok
  Let $X_i$ be iid with tail quantile function $U$ second-order regularly varying with indices
  $\gamma > 0$, $\rho \le 0$ and auxiliary function $A$. If $k \to \infty$, $k/n \to 0$ and
  $\sqrt k A(n/k) \to \lambda$, then
  \begin{align*}
    \sqrt k \, (\hat\gamma_{k,n} - \gamma) \; \Rightarrow \; N\big(\lambda/(1 - \rho), \gamma^2\big)
  \end{align*}
  [dHF, Theorem~3.2.5]. For exact Pareto samples, $\sqrt k(\hat\gamma - 1/\alpha) \Rightarrow
  N(0, \alpha^{-2})$ without the condition $k/n \to 0$.
\end{theorem}
\begin{proof}
  \leanok
  Write $X = U(e^E)$ with $E$ standard exponential. On the exponential scale the Hill estimator is
  $\gamma$ times the mean of $k$ iid exponential excesses plus a bias term
  $k^{-1}\sum_i (\log U(e^{E_{(i)}}) - \log U(e^{E_{(k+1)}}) - \gamma(E_{(i)} - E_{(k+1)}))$. The CLT
  handles the first part; the uniform second-order bound and $E_{(k+1)} - \log(n/k) \to 0$ show the
  bias is $A(n/k)/(1-\rho) + o_P(k^{-1/2})$.
\end{proof}

\begin{theorem}[Consistency of the moment estimator for all $\gamma$]
  \label{thm:moment-general}
  \uses{thm:mda-quantile, thm:extended-rv, thm:moment-consistency}
  \lean{ExtremeValueTheory.gpdMoment_one, ExtremeValueTheory.gpdMoment_two, ExtremeValueTheory.tendstoInMeasure_logMoment_of_isInMDAGEV,
    ExtremeValueTheory.tendstoInMeasure_momentEstimator_of_isInMDAGEV'}
  \leanok
  For $F$ in the domain of attraction of $H_\gamma$, $\gamma \in \bR$, with right endpoint $x_F > 0$,
  $k \to \infty$ and $k/n \to 0$: with $\gamma_- = \min(\gamma, 0)$ and $q_n = a(n/k)/U(n/k)$,
  $M^{(1)}/q_n \to 1/(1 - \gamma_-)$ and $M^{(2)}/q_n^2 \to 2/((1-\gamma_-)(1-2\gamma_-))$ in
  probability, and the moment estimator converges to $\gamma$ [dHF, Theorem~3.5.2].
\end{theorem}
\begin{proof}
  \leanok
  $\log U$ is extended regularly varying with index $\gamma_-$; a weak law at the deterministic
  threshold $U(n/k)$ via the layer-cake formula and Potter bounds, then sandwiching the random
  threshold $X_{(k+1)}$ with exceedance counts.
\end{proof}

\begin{theorem}[Weissman estimator under second-order regular variation]
  \label{thm:weissman-second-order}
  \uses{thm:weissman, thm:hill-normality, thm:second-order-bound}
  \lean{MeasureTheory.TendstoInDistribution.exists_measure_le_abs_le,
    ExtremeValueTheory.isBoundedInProbability_of_tendstoInDistribution, ExtremeValueTheory.isBoundedInProbability_hill,
    ExtremeValueTheory.IsSecondOrderRegularlyVarying.exists_bound_of_neg,
    ExtremeValueTheory.survivalQuantile_of_one_lt, ExtremeValueTheory.upperQuantile_eq_survivalQuantile,
    ExtremeValueTheory.IsSecondOrderRegularlyVarying.tendsto_div_mul_rpow, ExtremeValueTheory.tendsto_upperQuantile_div_of_secondOrder,
    ExtremeValueTheory.tendstoInMeasure_weissman_div_of_secondOrder,
    ExtremeValueTheory.tendstoInMeasure_tailProbEstimator_div_of_secondOrder, ExtremeValueTheory.tendstoInDistribution_weissman}
  \leanok
  Assume the tail quantile function is second-order regularly varying with $\gamma > 0$ and
  $\rho < 0$, $k \to \infty$, $k/n \to 0$ and $\sqrt k A(n/k) \to \lambda$. If $d_n = k/(n p_n) \ge 1$ and
  $\log d_n / \sqrt k \to 0$, then $\hat x_{p_n}/x_{p_n} \to 1$ in probability, and similarly for the
  tail-probability estimator. If moreover $d_n \to \infty$,
  \begin{align*}
    \frac{\sqrt k}{\log d_n} \Big( \frac{\hat x_{p_n}}{x_{p_n}} - 1 \Big)
      \; \Rightarrow \; N\big(\lambda/(1-\rho), \gamma^2\big)
  \end{align*}
  [dHF, \S4.3]. The case $\rho = 0$ is not covered.
\end{theorem}
\begin{proof}
  \leanok
  Convergence in distribution implies tightness, so the Hill CLT gives the rate. For $\rho < 0$ the
  uniform bound $|\log U(xt) - \log U(t) - \gamma \log x| \le C|A(t)|$ for all $x \ge 1$ makes the bias
  ratio tend to $1$. For normality, $\log(\hat x/x) = (\hat\gamma - \gamma)\log d_n + O_P(k^{-1/2})$ and the
  delta method for $\exp$.
\end{proof}

\section{The Hill estimator for dependent data}

\begin{theorem}[Hill consistency from a variance bound]
  \label{thm:hill-variance}
  \uses{thm:hill-consistency}
  \lean{ExtremeValueTheory.tailCount, ExtremeValueTheory.tendstoInMeasure_hill_of_variance, ExtremeValueTheory.tendsto_measure_abs_tailCount_div_sub,
    ExtremeValueTheory.logExcess_le_sum, ExtremeValueTheory.sum_le_logExcess}
  \leanok
  Let $X_i$ have identical marginals with $\bar F \in RV_{-\alpha}$, $k \to \infty$, $k/n \to 0$. If
  $\Var(\sum_{i<n} \mathbf 1\{X_i > u_n\}) / k_n^2 \to 0$ along every $u_n \to \infty$ with
  $n \bar F(u_n)/k_n \to c > 0$, then the Hill estimator converges to $1/\alpha$ in probability. No
  independence, stationarity or mixing is assumed.
\end{theorem}
\begin{proof}
  \leanok
  Chebyshev for the exceedance counts. The log-excess sum is squeezed deterministically between Riemann
  sums of exceedance counts on a geometric grid of thresholds; the remainder is controlled by Markov's
  inequality and Potter bounds, which use only the marginal law.
\end{proof}

\begin{theorem}[Hill consistency under m-dependence and strong mixing]
  \label{thm:hill-dependent}
  \uses{thm:hill-variance, def:mixing}
  \lean{ExtremeValueTheory.variance_tailCount_le_of_isMDependent, ExtremeValueTheory.tendstoInMeasure_hill_of_isMDependent,
    ExtremeValueTheory.variance_tailCount_le_of_isStationary, ExtremeValueTheory.tendstoInMeasure_hill_of_alphaMixing,
    ExtremeValueTheory.tendstoInMeasure_hill_of_alphaMixing_rate}
  \leanok
  The Hill estimator is consistent for $m$-dependent sequences with identical marginals (stationarity is
  not needed), and for stationary $\alpha$-mixing sequences with block length $r_n$ such that
  $\sum_{1 \le h \le r_n} \PR[X_h > u_n \mid X_0 > u_n] = o(k_n)$ and $(n/k_n)^2 \alpha(r_n + 1) \to 0$
  [Resnick--St\u{a}ric\u{a} 1995], [Rootz\'en--Leadbetter--de Haan 1998].
\end{theorem}
\begin{proof}
  \leanok
  $\Var \le (2m+1) n \bar F(u)$ under $m$-dependence. Under stationarity,
  $\Var \le n(\bar F(u) + 2\sum_{h \le r} \PR[X_0 > u, X_h > u] + 2 \sum_{h > r} \alpha(h))$.
\end{proof}

\begin{theorem}[Extreme quantiles at $\rho = 0$: bias cancellation]
  \label{thm:weissman-rho-zero}
  \uses{thm:weissman-second-order}
  \lean{ExtremeValueTheory.tendstoInMeasure_weissman_div_of_le, ExtremeValueTheory.tendstoInMeasure_weissman_div_of_antitoneOn,
    ExtremeValueTheory.tendstoInMeasure_tailProbEstimator_div_of_le, ExtremeValueTheory.tendstoInDistribution_weissman_of_tendsto_bias,
    ExtremeValueTheory.tendstoInDistribution_weissman_of_rho_eq_zero,
    ExtremeValueTheory.tendstoInDistribution_tailProbEstimator_of_tendsto_bias, ExtremeValueTheory.tendstoInDistribution_tailProbEstimator,
    ExtremeValueTheory.tendstoInDistribution_tailProbEstimator_of_rho_eq_zero}
  \leanok
  The Weissman and tail-probability results extend to $\rho \le 0$ given a majorant of $|A|$ on the
  extrapolation range (automatic when $|A|$ is eventually monotone). At $\rho = 0$, with $A$ asymptotically
  constant on that range, the Hill bias $+\lambda$ and the extrapolation bias $-\lambda$ cancel:
  $\sqrt k/\log d_n\,(\hat x/x_p - 1) \Rightarrow N(0, \gamma^2)$ and
  $\sqrt k/\log e_n\,(\hat p/p - 1) \Rightarrow N(0, \gamma^{-2})$. For $\rho < 0$ the tail-probability
  estimator satisfies $\sqrt k/\log e_n\,(\hat p/p - 1) \Rightarrow N(\lambda/((1-\rho)\gamma^2), \gamma^{-2})$
  [dHF, \S4.3--4.4].
\end{theorem}
\begin{proof}
  \leanok
  A uniform bound $|\log U(xt) - \log U(t) - \gamma \log x| \le C(1 + \log x)\,a$ when $|A| \le a$ on $[t, xt]$, and
  the delta method. A bound $C|A(t)|\log x$ fails in general for slowly varying $|A|$.
\end{proof}
